\section{Introduction}
\label{sec:introduction}

Differential cryptanalysis, introduced by Biham and Shamir [4], is one of
the most influential statistical attacks against symmetric-key primitives.
Its central insight is that the propagation of a chosen input difference
through an iterated cipher is far from uniform, and that pairs of input and
output differences that occur with high probability can be combined into
multi-round trails that distinguish the cipher from a random permutation or
recover key material. In the decades since its introduction, the technique
has been refined into a broad family of variants, including truncated,
higher-order, impossible, and boomerang differential attacks. All of them
share the same underlying object: the difference distribution table of the
cipher's nonlinear components, which records the frequency with which each
output difference is produced by each input difference.

The size of this table is also its main limitation. For an $n$-bit S-box,
the DDT is a $2^n \times 2^n$ integer matrix, and its quadratic blow-up in
$n$ rapidly exhausts memory budgets. An $8$-bit S-box already requires $64$
KiB; beyond $n = 10$, the dense matrix can no longer be materialized on
commodity machines. As cryptanalysts increasingly turn their attention to
larger S-boxes, ARX components, and word-oriented permutations, this
exponential cost forces a recurring compromise between the expressiveness of
differential analysis and the size of the primitives that can actually be
analyzed.

Compounding this storage problem is a more subtle access problem. Even when
a DDT fits in memory, most cryptanalytic tasks do not need every cell at
once: they need to extract the small subset of entries that satisfy some
structural property. Examples are pervasive in the literature, including
enumerating all impossible differentials of an S-box, listing every $(a,b)$
that achieves the differential uniformity, isolating truncated differentials
in which only a few input bits are active, identifying entries that satisfy
a linear relation imposed by the surrounding cipher, and filtering
transitions whose probability lies above a chosen threshold. In the standard
matrix view, each of these queries is solved by an ad-hoc procedure, written
anew for every analysis, and each shares the same unavoidable lower bound:
a full or partial sweep of the table.

In this paper, we argue that the right abstraction for the DDT is not a
matrix but a formal language. By recursively partitioning the table into
quadrants, every cell is identified with a unique word over a quaternary
alphabet, and the set of words that correspond to non-zero entries forms a
regular language $L(D)$. This language is recognized by a finite
automaton whose terminal states carry the differential multiplicities, and
whose structure is canonical once a quaternary analogue of multi-terminal
binary decision diagram reduction is applied. The resulting representation
exploits both the general redundancy present in every DDT, such as the
abundance of zero entries, and the specific redundancy resulting from the
algebraic and combinatorial structure of typical S-boxes.

The compactness of this representation is, however, only one half of the
contribution. The other half, and the focus of this work, is the querying
capability that the language-theoretic view unlocks. Once the DDT is
encoded as a regular language, any subset of entries definable by a regular
condition over the alphabet $\Sigma = \{\texttt{0},\texttt{1},\texttt{2},\texttt{3}\}$ becomes a regular
language in its own right, and can be intersected with $L(D)$ using
standard automata-theoretic constructions. The disparate ad-hoc scans
listed above all collapse into a single uniform recipe: write the property
as a regular expression, compile it into a query automaton, intersect with
the DDT automaton, and enumerate the resulting language. Truncated
differentials, impossible differentials, fixed-bit differentials, entries
of extremal probability, and many other targets that have historically
required specialized code are recovered as one-line queries. Importantly,
this also lifts queries from a per-cell cost to a cost proportional to the
size of the (already compressed) intersection automaton, so the dense
matrix never has to be enumerated even implicitly.

The remainder of the paper is organized as follows. Section~\ref{sec:preliminaries} recalls the
necessary background on differential cryptanalysis and binary decision
diagrams. Section~\ref{sec:automaton} describes the construction of the DDT automaton and
its reduction. Section~\ref{sec:results} reports memory measurements on a representative
collection of S-boxes. Section~\ref{sec:applications} then turns to applications, with
Section~\ref{sec:queries} developing the query framework that motivates this work.
Section~\ref{sec:conclusion} concludes and discusses directions for future research.
